\begin{table*}[t!]
\centering
%\small
%\rowcolors{2}{Para}{white}
\arrayrulecolor{black}
\resizebox{\textwidth}{!}{
\begin{tabular}{
>{\raggedright\arraybackslash}m{3cm}
>{\raggedright\arraybackslash}m{2.5cm}
>{\raggedright\arraybackslash}m{1cm}
>{\raggedright\arraybackslash}m{1cm}
>{\raggedright\arraybackslash}m{1cm}
>{\raggedright\arraybackslash}m{1cm}
>{\raggedright\arraybackslash}m{3cm}
>{\raggedright\arraybackslash}m{6cm}
>{\raggedright\arraybackslash}m{1cm}
}
\toprule
 \textbf{Method} 
& \textbf{Type} 
& \textbf{PR} 
& \textbf{TF} 
& \textbf{BUS} 
& \textbf{SUO} 
& \textbf{LMs} 
& \textbf{Main Advancement}
& \textbf{Year} \\
\toprule

\textbf{ULD} \newline\cite{ji2024reversing}
& additional parameters
& \usym{2713}
& \usym{2717}
& \usym{2713}
& \usym{2717}
& llama2-chat-7b, mistral-7b-instruct
& Derives an unlearned LLM by computing the \textbf{logit difference} between the target and assistant models.
& 2024 \\

\textbf{ECO} \newline\cite{liu2024large}
& prompt
& \usym{2713}
& \usym{2717}
& \usym{2713}
& \usym{2717}
& 68 LLMs ranging from 0.5B to 236B
& Performs unlearning by \textbf{corrupting prompt embeddings} detected by a classifier, without altering model weights.
& 2024 \\

\textbf{WAGLE} \newline \cite{jia2024wagle}
& locating-then-unlearning
& \usym{2717}
& \usym{2717}
& \usym{2713}
& \usym{2717}
& llama2-7b-chat, zephyr-7b-beta, llama2-7b
& Uses bi-level optimization to compute weight attribution scores for \textbf{selective fine-tuning} to achieve efficient, modular unlearning.
& 2024 \\

\textbf{SOUL}\newline\cite{jia2024soul}
& training objective
& \usym{2717}
& \usym{2717}
& \usym{2713}
& \usym{2713}
& opt-1.3b, llama2-7b
& Leverages a \textbf{second-order optimizer} for more effective LLM unlearning.
& 2024 \\

\textbf{SKU} \newline\cite{liu2024towards}
& training objective
& \usym{2717}
& \usym{2717}
& \usym{2713}
& \usym{2713}
& opt-2.7b, llama2-7b, llama2-13b
& Combines harmful knowledge learning with task vector negation in a two-stage framework for robust unlearning.
& 2024 \\

\textbf{EUL}\newline\cite{chen2023unlearn}
& additional parameters
& \usym{2713}
& \usym{2717}
& \usym{2713}
& \usym{2713}
& t5-base, t5-3b
& Introduces \textbf{unlearning layers} to forget specific data, supporting sequential unlearning through a \textbf{fusion mechanism} to merge multiple layers.
& 2023 \\

\textbf{ICUL} \newline\cite{pawelczyk2024context}
& prompt
& \usym{2713}
& \usym{2713}
& -
& -
& bloom-560m, bloom-1.1b, bloom-3b, llama2-7b
& First method to leverage \textbf{in-context learning (ICL)} for unlearning in language models.
& 2023 \\

\textbf{GA+Mismatch} \newline\cite{yao2024large}
& training objective
& \usym{2717}
& \usym{2717}
& \usym{2713}
& \usym{2717}
& opt-1.3b, opt-2.7b, llama2-7b
& Pioneered LLM unlearning by blending \textbf{forgetting, random mismatch, and KL-based preservation} objectives.
& 2023 \\

\textbf{KGA} \newline\cite{wang2023kga}
& training objective
& \usym{2717}
& \usym{2717}
& \usym{2713}
& \usym{2717}
& bart-base, distil-bert, lstm
& Simulates forgetting by aligning knowledge gaps between retain and forget models via \textbf{distributional divergence minimization}.
& 2023 \\

\textbf{DEPN}\newline\cite{wu2023depn}
& locating-then-unlearning
& \usym{2713}
& \usym{2713}
& \usym{2713}
& \usym{2717}
& bert-base
& Detects and disables privacy-related neurons to reduce sensitive data leakage in language models.
& 2023 \\

\bottomrule
\end{tabular}
}
\caption{Overview of methods for \textbf{parametric memory optimization in unlearning}. "PR" (Parametric Reserving) indicates whether the method avoids direct modification of internal weights. "TF" (Training-Free) shows if the method works without iterative optimization. "BUS" (Batch Unlearning Support) marks support for multiple edits simultaneously. "SUO" (Sequential Unlearning Optimization) indicates sequential unlearning capabilities. "LMs" lists language models used for experiments.}
\label{tab:methods-unlearning}
\end{table*}
